
\begin{table*}[t]
\centering
\caption{Positioning of the present work relative to recent bearing fault diagnosis studies.}
\label{tab:literature_positioning}
\small
\setlength{\tabcolsep}{5pt}
\renewcommand{\arraystretch}{1.18}
\begin{tabularx}{\textwidth}{>{\raggedright\arraybackslash}p{3.2cm}
>{\raggedright\arraybackslash}X
>{\raggedright\arraybackslash}X}
\toprule
Study direction & Common emphasis and limitation & Position of this study \\
\midrule
Deep supervised diagnosis and CNN/Transformer models & Emphasizes high classification accuracy under known fault classes, but usually requires labeled faults and can rely on condition-overlapping splits. & Studies normal-only anomaly detection and reports strict generalization protocols. \\
Autoencoder and VAE novelty detection & Uses reconstruction error as an anomaly score, while threshold choice and alarm behavior are often treated as secondary details. & Treats calibration and health decision rules as primary workflow components. \\
Self-supervised bearing representation learning & Learns transferable features from unlabeled signals, but is often evaluated as downstream classification. & Keeps the deployment task as abnormality screening without fault labels in training. \\
Domain adaptation and open-set diagnosis & Targets load, speed and dataset shift, but may assume target-domain adaptation data or source labels. & Compares holdout protocols and calibration behavior across CWRU and Paderborn. \\
Fuzzy health assessment & Provides interpretable condition categories, but is often separated from modern anomaly scores. & Connects fuzzy membership functions to calibrated anomaly scores. \\
\bottomrule
\end{tabularx}
\end{table*}
